\section{Canonical centre-ring scalars on the original normal graded layers}

Let $R$ be the actual original ambient holomorphic local ring, $I$ the
whole original ideal and $J$ the whole normal ideal.  The already constructed
layer $L_q$ is the image of $J^q$ in $R/(I+J^{q+1})$; its addition and its
intrinsic graded product are retained.  The following scalar descent is a
local algebraic ingredient for the Hilbert--Samuel equisingularity route in
Bierstone--Milman, subsection ``Use of the Hilbert-Samuel
function''.  That source passage describes the geometric goal; the explicit
quotient-scalar construction below is an additional algebraic argument.
It does not establish a normal-flat sheaf, centre containment, full ambient
Hilbert--Samuel comparison or descent under blowing up.  No action of $R/J$
on the full higher truncation $R/(I+J^s)$ is asserted.

\begin{lemma}[Actual ambient multiplication on the whole original layer]
\label{hir:relative-normal-actual-layer-ambient-scalars}
\source{bierstone-milman08-functoriality}
\uses{hir:relative-normal-graded-literal-layers,hir:relative-normal-intrinsic-layer-multiplication}
\lean{Hironaka.HironakaRelativeNormalCentreScalars.originalNormalGradedPieceAmbientModule,Hironaka.HironakaRelativeNormalCentreScalars.idealPowerLayerAmbient_smul_class}
\leanok
For any $q\geq0$, ordinary quotient multiplication defines a genuine
$R$-module on the existing literal layer $L_q$.  For every actual
$g\in J^q$ and $r\in R$, it is given by
\[
 r\,[g]_q=[rg]_q.
\]
The additive group is the original one; the action does not use a
polynomial quotient identification or a chosen tangent splitting.
\end{lemma}
\begin{proof}
Every element of $L_q$ has a representative $g\in J^q$ by its image
construction.  Since $J^q$ is an ideal, $rg\in J^q$.  Thus the actual
quotient product of the class of $r$ and this element is again in the
same image.  Multiplication in the original quotient is already
well-defined: if $g-g'\in I+J^{q+1}$, then $r(g-g')$ is in this same ideal.
Associativity, the scalar unit law and distributivity follow from the
corresponding equalities in that quotient, since the inclusion of its
image subtype is injective.  Its actual quotient class formula gives
the displayed equality for every representative, not only a chosen one.
\end{proof}

\begin{lemma}[Internally proved annihilation and canonical centre scalars]
\label{hir:relative-normal-actual-layer-centre-scalars}
\source{bierstone-milman08-functoriality}
\uses{hir:relative-normal-actual-layer-ambient-scalars}
\lean{Hironaka.HironakaRelativeNormalCentreScalars.originalNormalGradedPiece_isTorsionBy_normalIdeal,Hironaka.HironakaRelativeNormalCentreScalars.originalNormalGradedPieceCentreModule,Hironaka.HironakaRelativeNormalCentreScalars.originalNormalGradedPieceCentre_mk_smul_class}
\leanok
The whole original ideal $J$ annihilates $L_q$.  Consequently the actual
ambient action descends canonically to a module over the actual centre
quotient $R/J$, with the formula
\[
 (r\bmod J)\,[g]_q=[rg]_q\qquad(g\in J^q).
\]
No annihilator or desired quotient-action hypothesis is supplied.
\end{lemma}
\begin{proof}
For $u\in J$ and $g\in J^q$, their actual product belongs to
$J J^q=J^{q+1}$, hence has zero class in $R/(I+J^{q+1})$.
The actual power-class surjection therefore proves $u x=0$ for EVERY
$x\in L_q$.  If $r-r'\in J$, distributivity gives
$rx-r'x=(r-r')x=0$.  Thus the ambient scalar additive homomorphism factors
through the genuine quotient $R/J$.  On a quotient representative its
value is the original ambient multiplication.  The original ambient
module laws descend along the surjection $R\to R/J$, proving all module
laws and the formula.  This also proves uniqueness among quotient scalar
actions whose scalar representatives have the stated actual action.
\end{proof}

\begin{lemma}[Canonical scalars on the original whole direct sum]
\label{hir:relative-normal-actual-associated-centre-scalars}
\source{bierstone-milman08-functoriality}
\uses{hir:relative-normal-actual-layer-centre-scalars,hir:relative-normal-intrinsic-graded-algebra}
\lean{Hironaka.HironakaRelativeNormalCentreScalars.originalNormalAssociatedGradedAmbientModule,Hironaka.HironakaRelativeNormalCentreScalars.originalNormalAssociatedGraded_isTorsionBy_normalIdeal,Hironaka.HironakaRelativeNormalCentreScalars.originalNormalAssociatedGradedCentreModule,Hironaka.HironakaRelativeNormalCentreScalars.originalNormalAssociatedGradedCentre_smul_of}
\leanok
The existing literal direct sum $G=\bigoplus_{q\geq0}L_q$ inherits the
actual componentwise ambient action and canonical $R/J$ action.  For
every degree inclusion $\iota_q$, actual scalar $r$ and actual $x\in L_q$,
\[
 (r\bmod J)\,\iota_q(x)=\iota_q((r\bmod J)x).
\]
Every component of this action is the original quotient class product,
and $J$ annihilates the entire direct sum, all degrees simultaneously.
\end{lemma}
\begin{proof}
Define the ambient action componentwise; the support stays finite because
scalar multiplication preserves zero.  The laws hold in every actual
component and hence in the direct sum.  For $u\in J$, its action on every
component is zero by the preceding lemma, so its action on the entire
direct sum is zero.  Define the quotient action componentwise from the already canonically
descended layer actions.  Its laws hold in each component.  Representing
a quotient scalar by $r\in R$ identifies this action with the actual
ambient action componentwise, so it is exactly its quotient descent.
The ambient action commutes with each degree inclusion, as can be
checked in every component.  The two scalar representative formulas
then give the displayed inclusion identity.  Thus descent of the whole sum agrees with the
componentwise already descended layer actions.
\end{proof}

\begin{lemma}[Comparison after canonical descent with the actual plane trace]
\label{hir:relative-normal-actual-centre-trace-scalar-dictionary}
\source{bierstone-milman08-functoriality}
\uses{hir:relative-normal-actual-layer-centre-scalars,hir:relative-normal-actual-associated-centre-scalars,hir:relative-normal-entire-plane-kernel,hir:relative-normal-whole-hadamard}
\lean{Hironaka.HironakaRelativeNormalCentreScalars.planeQuotientAlgEquiv_mk,Hironaka.HironakaRelativeNormalCentreScalars.originalNormalGradedPieceCentre_smul_eq_planeQuotient,Hironaka.HironakaRelativeNormalCentreScalars.originalNormalGradedPieceCentre_mk_smul_eq_trace,Hironaka.HironakaRelativeNormalCentreScalars.originalNormalGradedPieceCentre_tangent_smul,Hironaka.HironakaRelativeNormalCentreScalars.originalNormalAssociatedGradedCentre_smul_eq_planeQuotient,Hironaka.HironakaRelativeNormalCentreScalars.originalNormalAssociatedGradedCentre_mk_smul_eq_trace,Hironaka.HironakaRelativeNormalCentreScalars.originalNormalAssociatedGradedCentre_tangent_smul}
\leanok
In original split coordinates, put $B=\mathcal O_{\mathbb C^d,0}$ and
$R=\mathcal O_{\mathbb C^{d+t},0}$.  The actual plane trace induces the
already proved algebra isomorphism $e:R/J\simeq B$.  For EVERY actual
quotient scalar $c$, the newly descended action on $L_q$ and $G$ agrees
with the existing tangent $B$-action of $e(c)$.  In particular
\[
 (r\bmod J)x=(\operatorname{trace} r)x,
 \qquad
 (\operatorname{pullback} b\bmod J)x=bx.
\]
The canonical centre action is constructed before this comparison.
\end{lemma}
\begin{proof}
The actual relative Hadamard decomposition says
$r-\operatorname{pullback}(\operatorname{trace}r)\in J$.
Its action on each actual layer is therefore zero.  The actual ambient
quotient multiplication by the pulled back $b$ is precisely the old
$tangent$ scalar action, both on every quotient numerator and on the
same image subtype.  Thus the first scalar formula holds for every
$r$ and every actual layer element.  Every quotient scalar has an
actual representative, and the true plane quotient isomorphism takes
this representative to its actual trace, so the formula holds for every
$c\in R/J$.  Trace of tangent pullback is the identity, giving the second
formula.  Componentwise verification gives the same two equalities on
the entire direct sum, with the original degree inclusions retained.
\end{proof}

Intrinsic free or flat modules over an actual smooth-centre stalk ring,
chart independence, a coherent sheaf/stalk identification and genuine
centre inclusion remain subsequent geometric producers.  The previous
same-open tangent ranks and fibre Hilbert--Samuel sequences do not by
themselves supply those conclusions.
